\documentclass[11pt,a4paper]{article}
\usepackage[T1]{fontenc}
\usepackage[utf8]{inputenc}
\usepackage{lmodern,amsmath,amssymb,array,longtable,booktabs}
\usepackage[margin=22mm]{geometry}
\usepackage{fvextra,xurl}
\usepackage[unicode,hidelinks]{hyperref}
\setlength{\emergencystretch}{3em}
\title{Exact identification of the AT extraction hash}
\author{}
\date{13 September 2026}
\begin{document}
\maketitle
This additive record identifies the processing frame of a retained metadata
hash. The entire accepted AT proof source and reader, including the original
preamble and all 23 equation tags, remain unchanged.

\section{The exact identified value}
The value in the sealed
\nolinkurl{evidence/BODY_EXTRACTION.json} field
\nolinkurl{pre_layout_reader_body_sha256} is
\nolinkurl{42b97470bfe0f67e0a0247e5cefcb95609c487b8ff20c275457aeebc02b7b706}.
It identifies the complete final \emph{post-layout} reader body in LF
encoding, with 17149 bytes. Replacing each CRLF in the 17558-byte accepted
reader by LF gives this value exactly. The retained field name and its
explanatory note name the wrong processing stage. The four exact frames are:

\begin{longtable}{@{}p{.23\linewidth}r p{.57\linewidth}@{}}
\toprule Complete body frame & Bytes & SHA-256 \\
\midrule
pre layout LF & 17113 & \nolinkurl{5ccb346f20b37ea30a40491a9c6ec84e7121869816efbb081007ae8f64cc9116} \\
pre layout CRLF & 17520 & \nolinkurl{3ed5d0dbedbcaffdaa8949a4dc6108e99546acd3023c32c4785388860a5d1815} \\
post layout LF & 17149 & \nolinkurl{42b97470bfe0f67e0a0247e5cefcb95609c487b8ff20c275457aeebc02b7b706} \\
post layout CRLF & 17558 & \nolinkurl{9fa19a509edf91f5aa0a9cb4762485449196d0b1ff521f2445f791fa3290cc75} \\
\bottomrule
\end{longtable}

\section{Complete forward and inverse reconstruction}
Read the pinned original source, extraction record and ten-entry layout
record. The exact original body begins with the 13-byte sequence whose
base64 representation is \nolinkurl{ClxtYWtldGl0bGUKCg==}.
Remove exactly that title prefix; no proof prose is removed.
The remaining complete LF body has 17113 bytes. Replacing each LF by CRLF
gives its 17520-byte CRLF representation.

Apply all ten ordered old/new pairs from the sealed layout map. In the
CRLF representation, convert every LF in both sides of each pair to CRLF.
Each old payload occurs exactly once at its recorded byte span. The final
LF and CRLF results are the complete final bodies identified above.
The full payloads, source spans and before/after hashes for every step are
retained in \nolinkurl{CLARIFICATION.json}.

For the inverse, start with the accepted CRLF reader and reverse all ten
pairs in reverse order, replacing each unique new payload by the complete
old payload. This recovers the entire pre-layout CRLF body. Replace each
CRLF by LF and restore the exact removed title prefix. Then concatenate
the original preamble, the literal begin-document boundary, the restored
complete body, the literal end-document boundary and the original
postamble. The result is every byte of the pinned original source,
whose SHA-256 is
\nolinkurl{24cf7ce5be6a43cc656641b516185df10b2e4d7824600ece657120d1c549deae}.
The JSON record gives exact zero-based byte intervals and hashes for all
five original-file components, and retains the complete original preamble
and postamble literally. The source itself retains every proof paragraph.

\section{Complete byte-verification program}
The following standard-library program checks the five sealed input pins,
the exact source spans and every intermediate hash in both representations,
the ten reverse transformations and all original-file components. It
writes no source file, imports no mathematical program and performs no
document build. Invoke it with \nolinkurl{--artifact-root} naming the
sealed artifact directory. Its complete JSON recipe accompanies this file.
\VerbatimInput[fontsize=\footnotesize,breaklines=true,breakanywhere=true]{verify_reconstruction.py}
\end{document}
